%\pagebreak
\section{Compound Directive Names}
\label{sec:compound}
\index{directive syntax!compound directive names}
\index{compound construct names}
\index{directive-name-modifier}

\index{compound construct@compound construct}
\index{compound construct@compound construct!combined construct@combined construct}
\index{compound construct@compound construct!composite construct@composite construct}
\index{compound directive names@compound directive names}
\index{compound directive names@compound directive names!leaf directive names@leaf directive name}
\index{compound directive names@compound directive names!constituent directive names@constituent directive name}

\index{construct properties@construct properties}
\index{construct properties@construct properties!parallelism-generating@parallelism-generating}
\index{construct properties@construct properties!thread-selecting@thread-selecting}
\index{construct properties@construct properties!partitioned@partitioned}

\index{clause properties@clause properties}
\index{clause properties@clause properties!all-constituents@all-constituents}
\index{clause properties@clause properties!outermost-leaf@outermost-leaf}
\index{clause properties@clause properties!innermost-leaf@innermost-leaf}
\index{clause properties@clause properties!all-privatizing@all-privatizing}
\index{clause properties@clause properties!once-for-all-constituents@once-for-all-constituents}

%%%%%%%%%%%%%%%%%

OpenMP directives can be classified as either \emph{leaf directives} or
\emph{compound directives}. Leaf directives are directives for which the
semantics are not defined in terms of other directives, while compound
directives are directives for which the semantics are defined in terms of
other directives, ultimately including the semantics of two or more leaf
directives. Compound directives may be further classified as either
\emph{combined directives} or \emph{composite directives}.

A combined directive is always expressed as a concatenation of two directive names,
\plc{directive-name-A} and \plc{directive-name-B}.
The first left-most directive name that appears in the directive is
\plc{directive-name-A},  and the remainder of the combined directive
name is \plc{directive-name-B}.  The semantics of a combined directive is
as if a construct formed by \plc{directive-name-A} immediately encloses
a construct formed by \plc{directive-name-B}.  An example of a
combined directive is \kcode{target teams loop}, which has the behavior of a
\kcode{target} construct that immediately encloses a \kcode{teams loop}
construct, and where the leaf directives are \kcode{target}, \kcode{teams}, and
\kcode{loop}.

A composite directive is any compound directive that is not combined. In
particular, it's not the case that the directive name can be split into a
\plc{directive-name-A} and \plc{directive-name-B}, such that the semantics are
a construct formed by \plc{directive-name-A} immediately enclosing a construct
formed by \plc{directive-name-B}. Usually, composite directives are also
expressed as a concatenation of two or more directive names, each explicitly
naming one of its leaf directives, where the first leaf directive name is
\plc{directive-name-A} and the remaining leaf directive names constitute
\plc{directive-name-B}. In general, such composite directives are
loop-associated, and the semantics are such that \plc{directive-name-A} forms a
loop-associated construct for which the loop body is a loop-associated
construct formed by \plc{directive-name-B}. An example of such a composite
directive is \kcode{distribute simd}, which has the behavior of a
\kcode{distribute} construct for which the loop body is a \kcode{simd}
construct, and where the leaf directives are \kcode{distribute} and
\kcode{simd}. Note that this is different than the \kcode{distribute} construct
immediately enclosing a \kcode{simd} construct, which in fact would not be
conforming OpenMP code. There are also composite directives that are not
expressed as a concatenation of two or more directive names, but nevertheless
include the semantics of two or more leaf directives. The \kcode{target_data}
directive is such a directive, entailing the semantics of a
\kcode{target_enter_data} construct, a \kcode{task} construct, and a
\kcode{target_exit_data} construct that are executed in sequence.

All OpenMP directives have one or more constituent directives.
If the OpenMP directive is not a compound directive (i.e., it is a leaf
directive) then it has only one constituent directive -- itself. Otherwise, for
a compound directive, the constituent directives include each of its leaf
directives and any directive that can be formed from a concatenation of
consecutive directive names that appear in the compound directive name. For
example, the constituent directives of a \kcode{target teams loop} directive
are: \kcode{target}, \kcode{teams}, \kcode{loop}, \kcode{target teams},
\kcode{teams loop}, and finally \kcode{target teams loop}. As another example,
the constituent directives of \kcode{target_data parallel loop} are:
\kcode{target_enter_data}, \kcode{task}, \kcode{target_exit_data},
\kcode{parallel}, \kcode{loop}, \kcode{target_data}, \kcode{target_data
parallel}, \kcode{parallel loop}, and finally \kcode{target_data parallel
loop}.

By default, clauses specified on a compound directive apply to the
individual leaf directives according to certain properties they may have:
\begin{itemize}
\item \plc{innermost-leaf} property: The clause only applies to the innermost
      resulting leaf construct that accepts the clause, though it may yield
      additional semantics for other leaf constructs. Examples:
      \kcode{private} and \kcode{linear} clauses.
\item \plc{outermost-leaf} property: The clause only applies to the outermost
      resulting leaf construct that accepts the clause. Examples:
      \kcode{nowait} and \kcode{nogroup} clauses.
\item \plc{all-privatizing} property: The clause only applies to those
      resulting leaf constructs to which a privatization clause is applied.
      Example: \kcode{allocate} clause.
\item \plc{once-for-all-constituents} property: The clause applies once for the
      most all-encompassing constituent construct. Example: \kcode{collapse}
      clause, which for a \kcode{target parallel for simd} directive would
      have the effect of collapsing the associated loop nest once for the
      \kcode{parallel for simd} constituent construct.
\item \plc{all-constituents} property: The clause applies to all
      resulting leaf constructs that accept the clause. This is the default
      property for a clause that does not have one of the above properties.
\end{itemize}

The OpenMP specification defines a grammar and various restrictions for
compound directive names that are composed of two or more explicitly named
constituent directive names. The grammar allows for a handleful of such
directive names that are composite, and additionally permits two directive
names to be concatenated into a combined directive name according to the
following rules:

\begin{itemize}
\item the first named directive (\plc{directive-name-A}) consists of only one
      explicitly named constituent directive that is either
      parallelism-generating or thread-selecting;
\item the second named directive (\plc{directive-name-B}) is the name of a
      parallelism-generating, thread-selecting, or partitioned directive, or any
      compound directive for which its \plc{directive-name-A} is such a
      directive.
\end{itemize}

For example, consider the directive \kcode{masked target teams distribute
simd}, which conforms to the above rules as shown in the following diagram:

\begin{Verbatim}[fontsize=\footnotesize]

    masked target teams distribute simd (combined)
          ______________|_____________
         |                            |
   A: masked       B: target teams distribute simd (combined)
(thread-selecting)        ________|__________________
                         |                           |
                     A: target            B: teams distribute simd (combined)
                 (parallelism-generating)  __________|____________
                                          |                       |
                                       A: teams         B: distribute simd (composite)
                               (parallelism-generating)    _______|_______
                                                          |               |
                                                     A: distribute      B: simd
                                                     (partitioned)

\end{Verbatim}

In contrast, the directive \kcode{parallel atomic} would not be conforming,
because \kcode{atomic} is not a permitted \plc{directive-name-B} according to
the above rules. For more details on the grammar and associated restrictions,
please refer to the \docref{Compound Directive Names} section of the OpenMP
specification.

The following example attempts to make use of various invalid compound
directives that are supported by the grammar but violate one of the
associated restrictions.

The \kcode{parallel target parallel} directive is invalid due to a restriction
that a compound directive cannot have multiple leaf directives with the same
name, such as repeated \kcode{parallel}. This avoids ambiguity when specifying a
particular leaf directive to which a clause should be applied with the
\plc{directive-name-modifier}. The \kcode{teams target} directive is invalid
due a restriction that a combined directive name must specify an inner
construct formed by \plc{directive-name-B} that may be immediately nested
inside an outer construct formed by \plc{directive-name-A}. Since a
\kcode{target} construct may not be nested inside a \kcode{teams} construct,
the combined directive is disallowed. The \kcode{teams sections} and
\kcode{teams masked} directives are both invalid due to a restriction that
worksharing and thread-selecting directive names may only combined with a
\kcode{parallel} \plc{directive-name-A}. Generally speaking, any other
\plc{directive-name-A} would not form a useful combined consruct. The
\kcode{target_data parallel target} directive is invalid due to a restriction
that at most one constituent construct may be map-entering (or map-exiting),
and here the \kcode{target_data} and \kcode{target} constituent constructs are
both map-entering and map-exiting. Finally, the \kcode{task taskloop} directive
is invalid due to a restriction that two task-generating constructs may not be
combined if they generate explicit tasks that bind to the same parallel region.

\cexample[6.0]{invalid_compound_names}{1}
\ffreeexample[6.0]{invalid_compound_names}{1}

In the following example, each loop-nest-associated directive is a compound
directive that applies different parallelization semantics according to its
constituent directives.

In the first construct of the \ucode{component_names} procedure, the directive
\kcode{parallel for} (or \kcode{parallel do}) forms a combined construct that
is equivalent to a \kcode{parallel} construct (a parallelism-generating
construct) that immediately encloses a \kcode{for} (or \kcode{do}) construct
(a partitioned construct).

In the second construct of the procedure, the directive \kcode{parallel for
simd} (or \kcode{parallel do simd}) forms a combined construct that is
equivalent to a \kcode{parallel} construct (a parallelism-generating
construct) that immediately encloses a \kcode{for simd} (or \kcode{do simd})
composite construct (where the first named leaf directive is partitioned).

In the third construct of the procedure, the directive \kcode{parallel masked
taskloop} forms a combined construct that is equivalent to a \kcode{parallel}
construct (a parallelism-generating construct) that immediately encloses
a \kcode{masked taskloop} combined construct. That combined construct, in turn,
is equivalent to a \kcode{masked} construct (a thread-selecting
construct) that immediately encloses a \kcode{taskloop} construct (a
parallelism-generating construct).

For the fourth construct of the procedure, the directive \kcode{target teams
distribute parallel for} (or \kcode{target teams distribute parallel do}) forms
a combined construct that is equivalent to a \kcode{target} construct (a
parallelism-generating construct) that immediately encloses a \kcode{teams
distribute parallel for} (or \kcode{teams distribute parallel do}) combined
construct. That combined construct, in turn, is equivalent to a \kcode{teams}
construct (a parallelism-generating construct) that immediately encloses
a \kcode{distribute parallel for} (or \kcode{distribute parallel do}) composite
construct (where the first named leaf directive is partitioned).

\cexample[6.0]{compound_names}{1}
\ffreeexample[6.0]{compound_names}{1}

In general, OpenMP permits the clause properties described earlier in this
section to be overridden for any clause with the use of a
\plc{directive-name-modifier} (prior to OpenMP 6.0, this modifier was only
permitted in the \kcode{if} clause). The modifier can name any constituent
directive of the directive on which it appears. For example, a
\kcode{private(parallel: x)} clause on a \kcode{parallel loop} directive would
apply the \kcode{private} clause to the \kcode{parallel} directive rather than
the default behavior of applying it to the \kcode{loop} directive due to the
innermost-leaf property.

In the next example, the first three constructs of the
\ucode{conditional_parallel} procedure apply an \kcode{if} clause to the
\kcode{parallel} construct in three different ways: one without a
\plc{directive-name-modifier} and the other two with
\plc{directive-name-modifiers} (one specifying a leaf directive name and the
other specifying a constituent compound directive name). Without the modifier,
the \kcode{if} clause applies to all the constituent directives that accept the
\kcode{if} clause (only the \kcode{parallel} here). With the \kcode{parallel}
\plc{directive-name-modifier}, it applies explicitly to only the
\kcode{parallel} directive. With \kcode{parallel for} (or \kcode{parallel do})
\plc{directive-name-modifier}, it applies explicitly to the \kcode{parallel
for} (or \kcode{parallel do}) directive, which has the effect of applying only
to the \kcode{parallel} leaf directive.

In the last construct of the procedure, the \plc{directive-name-modifier} specifies
that the \kcode{if} clause should only apply to the \kcode{taskloop simd}
constituent directive. This has the effect of applying the \kcode{if} clause to
the \kcode{taskloop} and \kcode{simd} leaf directives.

\cexample[6.0]{directive_name_modifier}{1}
\ffreeexample[6.0]{directive_name_modifier}{1}
